\documentclass[11pt]{article}
\usepackage[margin=1in]{geometry}
\usepackage{amsmath,amssymb,amsthm,mathtools}
\usepackage[hidelinks]{hyperref}
\newtheorem{theorem}{Theorem}
\newtheorem{lemma}[theorem]{Lemma}
\newtheorem{corollary}[theorem]{Corollary}
\title{Pad\'e strip separation for adjacent nonconstant monomial supports}
\author{Internal algebraic proof and symbolic positivity certificates; no priority claim}
\date{17 November 2026}
\begin{document}
\maketitle
All moments use ordinary factorials. We study
$f(X)=X^a(A+BX)$, where $a$ is a nonnegative integer and $A,B$ are
complex, with a nonzero polynomial. This changes the support family:
when $a>0$ both monomials are nonconstant and the Gamma shape in the
$n$th moment is $an+1$, varying with $n$. The fixed-shape Appell
squarefreeness equivalence is not used.

\begin{theorem}\label{thm:family}
For every $a\in\{1,2,3,4,5,6\}$, $X^a(A+BX)$ satisfies Strong Factorial
at every starting index. The case $a=0$ is already covered by the
\href{ALL_UNIVARIATE_TWO_TERM_SFC.tex}{constant-plus-power theorem}.
More generally, for every pair of integers $n\ge2$, $a\ge1$ with
\[
                            a\ge2n-3,                 \tag{1}
\]
the two moments at orders $n-1,n$ cannot both vanish.
\end{theorem}
The retained \href{sources/rocks-2015-author-preprint.pdf}{Rocks 2015 author manuscript},
Proposition 3.14, already gives
the $a=0,1,2$ cases; Proposition 3.15 and Remarks 3.18--3.19 give partial
statements for general $a\ge3$. This is a scope comparison with the pinned
author draft, not a novelty claim or complete final-publication audit.
The general $a$ assertion, arbitrary monomial pairs and unrestricted
FC/SFC remain open beyond the conclusions proved here.

\section{Exact Pad\'e polynomials, proved from coefficients}
For integers $M,N\ge0$, put
\[
 P_{M,N}(z)=\sum_{j=0}^M
 \frac{(M+N-j)!M!}{(M+N)!(M-j)!j!}z^j,
\quad
 Q_{M,N}(z)=\sum_{j=0}^N
 \frac{(-1)^j(M+N-j)!N!}{(M+N)!(N-j)!j!}z^j.
                                                               \tag{2}
\]
Both constant terms are one, and their degrees are exactly $M,N$.
They satisfy
\[
                  e^zQ_{M,N}(z)-P_{M,N}(z)=O(z^{M+N+1}). \tag{3}
\]
For a direct coefficient proof, multiply $[z^j]e^zQ$ by
$(M+N)!j!/N!$. For $0\le j\le M+N$, the result is
\[
 \sum_{k=0}^{j}(-1)^k\binom jk(M+N-k)_{\underline M}.
\]
Terms with $k>N$ vanish, so extending the original sum to $j$ is valid.
The finite difference is zero for $j>M$, and otherwise equals
$(M)_{\underline j}(M+N-j)_{\underline{M-j}}$.
This is exactly the coefficient of $P$ in (2), proving (3) without
an imported Pad\'e normality theorem.

For the normalized factorial moment set
\[
 B_{n,a}(t)=\sum_{k=0}^n\binom nk(an+1)_k t^k.
\]
Then
\[
 \mathcal L(f^n)=A^n(an)!B_{n,a}(B/A),\qquad
 Q_{an,n}(z)=\frac{(-1)^n(an)!}{((a+1)n)!}
                  z^nB_{n,a}(-1/z).                 \tag{4}
\]
Thus, for nonzero $A,B$, the adjacent obstruction at terminal $n$ is
exactly a common zero of $Q_{an,n}$ and $Q_{a(n-1),n-1}$.
No zero at $z=0$ occurs. If either coefficient vanishes, $f$ is a monomial
and all its positive moments are nonzero.

\section{A sharp elementary strip for the denominators}
\begin{lemma}\label{lem:strip}
Every zero $\zeta$ of $Q_{M,N}$, for $N\ge1$, satisfies
\[
          M-N+2\le\operatorname{Re}\zeta\le M+N.       \tag{5}
\]
Both inequalities are strict when $N\ge2$.
\end{lemma}
\begin{proof}
Write $S=M+N$ and consider the $N$ by $N$ complex symmetric tridiagonal
matrix $T$ with
\[
 T_{jj}=S+2-2j,\qquad
 T_{j,j+1}=T_{j+1,j}=i\sqrt{j(S+1-j)}\quad(1\le j<N).
\]
All square roots are positive real. Its leading principal characteristic
polynomials $D_j$ satisfy
\[
 D_0=1,\quad D_1=z-S,\quad
 D_j=(z-S-2+2j)D_{j-1}+(j-1)(S+2-j)D_{j-2}.
\]
Substitution of
\[
 D_j(z)=\sum_{k=0}^j(-1)^{j-k}\frac{j!}{k!}
                      \binom{S-k}{j-k}z^k
\]
verifies the recurrence and its initial values. At $j=N$ it gives
\[
            Q_{M,N}(z)=\frac{(-1)^N M!}{(M+N)!}\det(zI-T).
\]
The Hermitian part $(T+T^*)/2$ is the real diagonal matrix with entries
$S,S-2,\ldots,S+2-2N$. For an eigenvector $u$,
$\operatorname{Re}\zeta=u^*((T+T^*)/2)u/(u^*u)$, proving (5).
When $N\ge2$, equality at an extreme would force $u$ to be supported
on the first or last coordinate only. The nonzero neighboring matrix
entry makes this impossible for an eigenvector. This proves strictness.
\end{proof}
Applying the strict lower bound to $Q_{an,n}$ and the upper bound to
$Q_{a(n-1),n-1}$ gives (1), since the two strip endpoints differ by
$a+3-2n$. Equality is harmless because the first lower bound is strict.

\section{The exact auxiliary polynomial}
Fix $a\ge1,n\ge2$ and put
\[
 \ell=(a+1)(n-1)+1,\qquad
 C_{a,n}=\frac{n!}{((a+1)n)!}
         \frac{(-1)^{n-1}(a(n-1))!}{((a+1)(n-1))!}\ne0.
\]
Equations (2)--(3) imply the polynomial identity
\[
 P_{an,n}Q_{a(n-1),n-1}-P_{a(n-1),n-1}Q_{an,n}
       =C_{a,n}z^\ell H_{a,n}(z).                    \tag{6}
\]
For $a\ge2$, $H$ is monic of degree $a-1$; for $a=1$, $H=2$.
Indeed the cross determinant vanishes to order at least $\ell$ by (3).
Its degree is at most $(a+1)n-1$, and its highest coefficients are obtained
from the two product expansions. Explicitly, if
$H=\sum_{k=0}^{a-1}h_kz^{a-1-k}$, then
\[
 h_k=\sum_{i+j=k}
 \frac{(n+1)_i(an)_{\underline i}}{i!}
 \frac{(-1)^j(a(n-1)+1)_j(n-1)_{\underline j}}{j!}
 +\mathbf1_{\{k=a-1\}}\frac{(an)_{\underline a}}n.     \tag{7}
\]
The final term is the contribution of the leading term of the second
product, which first appears at degree $\ell$. These computations prove
(6)--(7) for all parameters; small-index replay is only a bookkeeping check.

Any common denominator zero must be a zero of $H$, since $z\ne0$.
By Lemma~\ref{lem:strip}, a sufficient separation condition is
\[
 H_{a,n}(c+w)\text{ has all roots in }\operatorname{Re}w<0,
                    \qquad c=(a-1)n+2.               \tag{8}
\]
This is sufficient, not claimed necessary for the SFC assertion.
It reduces an $n$-dependent degree-$n$ obstruction to an explicit
polynomial of degree $a-1$.

For later exact verification, (7) gives
\[
 H_{2,n}=z+9n-3,
\quad H_{3,n}=z^2+2(4n-1)z+8(2n-1)(4n-1),
\]
\[
 \begin{split}
 H_{4,n}={}&z^3+(11n-3)z^2+3(5n-2)(5n-1)z\\
          &+5(5n-3)(5n-2)(5n-1),
 \end{split}
\]
\[
 \begin{split}
 H_{5,n}={}&z^4+2(7n-2)z^3+6(3n-1)(7n-2)z^2\\
 &+24(2n-1)(3n-1)(6n-1)z\\
 &+72(2n-1)(3n-2)(3n-1)(6n-1),
 \end{split}
\]
\[
 \begin{split}
 H_{6,n}={}&z^5+(17n-5)z^4+10(19n^2-12n+2)z^3\\
 &+2(7n-3)(7n-2)(17n-5)z^2\\
 &+5(7n-4)(7n-3)(7n-2)(7n-1)z\\
 &+7(7n-5)(7n-4)(7n-3)(7n-2)(7n-1).
 \end{split}
\]

\section{All-index shifted separation by dissipative matrices}
\begin{lemma}\label{lem:schwarz}
Let a real monic polynomial $p$ of degree $d\ge1$ be generated by
\[
 D_0=1,\quad D_1=w+b_0,\quad
 D_{j+1}=wD_j+b_jD_{j-1}\ (1\le j<d),\quad D_d=p,
                                                               \tag{9}
\]
with every $b_j>0$. Then all roots of $p$ have strictly negative real part.
\end{lemma}
\begin{proof}
Take the real tridiagonal matrix with diagonal $(-b_0,0,\ldots,0)$,
upper entries $-\sqrt{b_j}$ and lower entries $\sqrt{b_j}$.
Its characteristic polynomial is (9). Its Hermitian part is
$\operatorname{diag}(-b_0,0,\ldots,0)$, so every eigenvalue has nonpositive
real part. Equality forces the first eigenvector coordinate to vanish;
the first row and then successive rows force every coordinate to vanish,
a contradiction. For $d=1$ the only root is $-b_0$.
\end{proof}

For each $a=2,3,4,5,6$, set $p(w)=H_{a,n}(w+(a-1)n+2)$.
The accompanying \href{pade_shift_positive_certificates.json}{exact symbolic
certificate} gives rational functions $b_0(n),\ldots,b_{a-2}(n)$ with
both of these properties:
\begin{enumerate}
\item Substitution into the finite recurrence (9) gives $p$ identically
      as a polynomial in $w$ over $\mathbb Q(n)$.
\item On writing $u=n-2$, the numerator and denominator of each $b_j$
      are polynomials in $u$ with nonnegative integer coefficients and
      strictly positive constant coefficient.
\end{enumerate}
Thus every parameter is strictly positive at every real $n\ge2$, with
no pole or omitted integer, and Lemma~\ref{lem:schwarz} proves (8).
For example the first two certificates are
\[
 a=2:\ b_0=10n-1;\qquad
 a=3:\ b_0=2(6n+1),\ b_1=4(21n^2-7n+2).
\]
For $a=4$ they are
\[
 \begin{split}
 b_0&=20n+3,\\
 b_1&=4(596n^3+355n^2-19n+10)/(20n+3),\\
 b_2&=(976n^3-576n^2+247n-22)/(20n+3).
 \end{split}
\]
The higher two certificates have numerator/denominator degrees at most
$11$ and $9$, respectively; all their exact coefficient lists are retained
in the certificate file and verified by polynomial multiplication.
These finite identities are universal evidence in $n$, not a finite scan
of moment indices.
For $a=1$, the nonzero constant $H=2$ already excludes every common root.
Together with (4)--(8), this proves Theorem~\ref{thm:family}.

The certificate can be constructed from the parity parts of $p$:
set $E$ to its leading-degree parity part, $O$ to the other part divided
by its leading coefficient $b_0$. Repeatedly normalize the remainder
$E-wO$ to recover $b_1,b_2,\ldots$. The written certificates check the
reconstruction (9) independently, so a successful result does not rely
on an unproved assertion that this procedure always yields positive data.
For any supplied $a$ where both exact conditions hold, the same theorem
certifies all moment indices; positivity for every $a$ is not established.

\section{A small explicit matrix pencil for the uniform gate}
There is an independent necessary polynomial condition with especially
simple matrix entries. Define
\[
 R_d(t,z)=\sum_{k=0}^d(-1)^k\binom dk(t+1)_k z^{d-k},\qquad
 \sum_{d\ge0}R_d(t,z)\frac{u^d}{d!}
              =e^{zu}(1+u)^{-t-1}.
\]
For nonnegative integer $t$, $Q_{t,d}$ is a nonzero constant multiple of
$R_d(t,z)$. Direct differentiation of the generating function gives
\[
 (t+1)R_d(t+1,z)=(t+d+1+z)R_d(t,z)-zR_d(t-1,z),
\]
\[
 R_d(t,z)=zR_{d-1}(t-1,z)-(t+d)R_{d-1}(t,z).          \tag{10}
\]
Suppose $a\ge2$ and $z$ is a common denominator zero, and write
$M=an$, $l=a(n-1)$, $D=a-1$, $S_t=R_{n-1}(t,z)$.
Then $S_l=0$ and $zS_{M-1}=(M+n)S_M$.
The vector $v_j=S_{l+j}$, $1\le j\le D$, is nonzero: otherwise
$S_l=S_{l+1}=0$, and the first recurrence propagates zeros to every
larger integer $t$. That contradicts the nonzero leading coefficient
$(-1)^{n-1}$ of the degree-$n-1$ polynomial $S_t$ in $t$.

Eliminate $S_M$ with the boundary relation. Equations (10) become
\[
                    (\mathsf A+z\mathsf B)v=0,        \tag{11}
\]
where $\mathsf A$ is upper bidiagonal, $\mathsf B$ is lower bidiagonal,
and, for $1\le j\le D$,
\[
 \mathsf A_{jj}=l+j+n,\qquad
 \mathsf A_{j,j+1}=-(l+j+1),\qquad
 \mathsf B_{jj}=\begin{cases}1,&j<D,\\1/(a+1),&j=D,\end{cases}
 \quad \mathsf B_{j,j-1}=-1.
\]
Consequently every genuine adjacent obstruction must satisfy
$\det(\mathsf A+z\mathsf B)=0$.
This is a polynomial of degree $a-1$, since
$\det\mathsf B=1/(a+1)\ne0$.
For the verified cases $2\le a\le6$, direct polynomial identities give
$(a+1)\det(\mathsf A+z\mathsf B)=H_{a,n}(z)$; the necessary-condition
argument does not require asserting this identity for unbounded $a$.

A sufficient unbounded target, now an explicit finite-matrix inequality, is
\[
 \operatorname{Re}\lambda(-\mathsf B^{-1}\mathsf A)<(a-1)n+2
                  \quad(a\ge2,\ n\ge2).              \tag{12}
\]
The matrix dimension depends on $a$, not on the moment index $n$, and
$\mathsf A$ is affine in $n$. Equation (12) remains unproved here.
No conclusion for all $a$ is inferred from the finite symbolic certificates
or the additional exact terminal diagnostics.

\section{The genuine remaining boundary}
The blanket unshifted shortcuts fail:
$H_{3,2}=z^2+14z+168$ has no real roots, and
$H_{5,2}=z^4+24z^3+360z^2+3960z+47520$ has negative third Hurwitz
determinant $24\cdot360\cdot3960-3960^2-24^2\cdot47520$.
Neither refutes SFC or the shifted target (8).
The uniform assertion (8) for all $a\ge2$ remains unproved here.
After the registered ranges, a residual gate is $a\ge7$,
$n>(a+3)/2$, unless a supplied positive shifted certificate closes a
particular $a$. The generic exporter is a sufficient certificate test;
failure is not a counterexample or an irreducibility claim.

No general $X^a(A+BX^d)$ result for $d>1$, arbitrary monomial pairs,
other Gamma shapes, general Factorial or unrestricted Strong Factorial
is inferred. The overarching goal stays open. Independent replay,
formal verification and external review remain false. All proofs in
this note are algebraic; no zero-distribution inference or unproved
asymptotic cutoff enters.
\end{document}
